\documentclass{scrartcl}

\usepackage[utf8]{inputenc}

\usepackage{amsmath}
\usepackage{babel}

\usepackage{xcolor}
\usepackage{tikz}
\usetikzlibrary{arrows.meta,calc,decorations.markings,arrows,positioning,fit,intersections,patterns,scopes,datavisualization,backgrounds}
\usetikzlibrary{shapes.geometric,decorations.pathreplacing,shadings,decorations.text}

\usepackage{hyperref}
\usepackage[footnotes,definitionLists,smartEllipses,hybrid,pipeTables=true,shiftHeadings=1,tableCaptions=true]{markdown}
\usepackage[edges]{forest}

\usepackage{booktabs}
\usepackage{cleveref}

\usepackage{paralist}

\usepackage{pifont}
\newcommand{\jaja}{\ding{51}}
\newcommand{\nope}{\ding{55}}

\usepackage{tabularx}
\renewcommand\tabularxcolumn[1]{m{#1}}% for vertical centering text in X column
\newcolumntype{Y}{>{\centering\arraybackslash}X}

\title{Appendix B: IFC Files and Descriptions (DEPL Part)}
\author{Files submitted to IFC Infrastructure Extension Deployment project}
\date{document automatically generated on: \today}

\setkeys{Gin}{width=\linewidth}

\usepackage{ltablex}
\usepackage{tabularx}
\usepackage{booktabs}

\usepackage{datetime2}

\usepackage{fancyhdr}
\pagestyle{fancy}
\lhead{}
\chead{\includegraphics[width=0.4\textwidth]{bSI_Logo2021.png}}
\rhead{}
\setlength\headheight{45pt} 

\makeatletter
\def\markdownLaTeXReadAlignments#1{%
\advance\markdownLaTeXColumnCounter by 1\relax%
\if#1d%
\addto@hook\markdownLaTeXTableAlignment{p{0.5\textwidth}}%
\else%
\addto@hook\markdownLaTeXTableAlignment{#1}%
\fi\ifnum\markdownLaTeXColumnCounter<\markdownLaTeXColumnTotal\relax\else\expandafter\@gobble%
\fi\markdownLaTeXReadAlignments%
}
\makeatother

\newcommand{\pr}[1]{\href{https://github.com/bSI-InfraRoom/IFC-infra-unit-test/pull/#1}{\##1}}

\begin{document}

\maketitle
\addtocontents{toc}{~\hfill\textbf{Page}\par}

\section{Summary}
\label{sec:summary}

This is a summary document of all unit tests submitted during the IFC Infrastructure Extensions Deployment project until \DTMdisplaydate{2021}{7}{16}{}.

This Section is structured as follows:
\begin{compactitem}
  \item guidance on how to read the document in \cref{sec:how_to},
  \item a list of all unit tests, their status of approval, and source of data in \cref{sec:status},
  \item an overview of the covered entities by all approved unit tests in \cref{sec:coverage},
  \item an overview over the dependencies between individual unit tests in \cref{sec:schematics}, and
  \item a table of content in \cref{sec:content}.
\end{compactitem}
\Cref{sec:project_setup,sec:alignment,sec:linear_placement,sec:swept_geometry,sec:domain_specifics,sec:earthworks,sec:geodetic_survey,sec:combinations,sec:georeferencing_and_dtms,sec:drainage_system,sec:other,sec:spatial_structure} are automatically generated from the \texttt{readme} files of individual unit tests.
The grouping of individual unit tests under these sections is (almost) arbitrary and should only serve for orientation.

This document has been set up by \v{S}tefan Jaud, validation lead of the project.
All content has been submitted to the project by participating software vendors (see Appendix A of the delivery).
The originating documentation together with all mentioned files can be obtained
 from the project's official GitHub repository\footnote{\url{https://github.com/bSI-InfraRoom/IFC-infra-unit-test}}.

\subsection{How to read the document}
\label{sec:how_to}

Each unit test is within its own subsection with a link to the original data on the GitHub.
The included figures are \emph{floating} in this document, so they may appear at the top of the page,
 even though they are mentioned later in the text.
All figures are always flushed before the next section starts, so that there is no doubt of their assignment to the unit test.

\subsection{Status}
\label{sec:status}

\Cref{tab:unit_tests} lists all unit tests and presents their status:
\begin{compactitem}
  \item column \emph{Section} provides a link to the corresponding section in the report,
  \item column \emph{Unit test} provides the name of the unit test,
  \item column \emph{Export} names the producing software vendor (or the project team itself if it was produced by us),
  \item column \emph{Ck.} shows, if the unit test file has been checked (\tikz{\jaja}) by the project team or not (\tikz{\nope}),
  \item column \emph{Import} names another software solution that successfully imported the data (can be empty if it has not happened yet), and
  \item column \emph{SL} denotes unit tests stemming from a storyline (\tikz{\jaja}).
\end{compactitem}

\begin{tabularx}{\textwidth}{rYlclc}
  \caption{An overview of the status of the unit tests. See \cref{sec:status} for an explanation of the columns and entries.}
  \label{tab:unit_tests} \\
  \toprule
    Section     & Unit test    &  Export     & Ck.       & Import      & SL \\
  \midrule
  \endfirsthead
  \caption[]{\emph{(cont.)} An overview of the status of the unit tests. See \cref{sec:status} for an explanation of the columns and entries.} \\
  \toprule
    Section     & Unit test    &  Export     & Ck.       & Import      & SL \\
  \midrule
  \endhead
  \bottomrule
  \multicolumn{6}{r}{\footnotesize(continued)}
  \endfoot
  \bottomrule
  \endlastfoot
  % content
    \labelcref{sec:project_setup_1} & Setup 1 & project team & \jaja  & TUM OIP & \nope \\
    \labelcref{sec:project_setup_2} & Setup 2 & project team & \jaja  & TUM OIP & \nope \\
  \midrule
    \labelcref{sec:tin_1} & DTM 1 & 12d Solutions & \jaja  & TUM OIP & \nope \\
    \labelcref{sec:georeferencing_1} & Georef. 1 & project team & \jaja  & TUM OIP & \nope \\
    \labelcref{sec:georeferencing_3} & Georef. 2 & project team & \jaja  & TUM OIP & \nope \\
    \labelcref{sec:georeferencing_4} & Georef. 3 & project team & \jaja  & TUM OIP & \nope \\
    \labelcref{sec:georef_tin_1} & Georef.+DTM 1 & 12d Solutions & \jaja  & TUM OIP & \nope \\
    \labelcref{sec:georef_tin_2} & Georef.+DTM 2 & 12d Solutions & \jaja  & TUM OIP & \nope \\
    \labelcref{sec:georef_tin_3} & Georef.+DTM 3 & 12d Solutions & \jaja  & TUM OIP & \nope \\
    \labelcref{sec:georef_tin_4} & Georef.+DTM 4 & 12d Solutions & \jaja  & TUM OIP & \nope \\
  \midrule
    \labelcref{sec:align_1} & Alignment 1 & 12d Solutions & \jaja  & Geometry Gym & \nope \\
    \labelcref{sec:align_2} & Alignment 2 & 12d Solutions & \jaja  & Geometry Gym & \nope \\
    \labelcref{sec:align_3} & Alignment 3 & 12d Solutions & \jaja  & Geometry Gym & \nope \\
    \labelcref{sec:align_4} & Alignment 4 & 12d Solutions & \jaja  & Geometry Gym & \nope \\
    \labelcref{sec:align_5} & Alignment 5 & 12d Solutions & \jaja  & Geometry Gym & \nope \\
    \labelcref{sec:align_6} & Alignment 6 & 12d Solutions & \jaja  & Geometry Gym & \nope \\
    \labelcref{sec:align_7} & Alignment 7 & 12d Solutions & \jaja  & Geometry Gym & \nope \\
    \labelcref{sec:align_8} & Alignment 8 & 12d Solutions & \jaja  & Geometry Gym & \nope \\
    \labelcref{sec:align_9} & Alignment 9 & 12d Solutions & \nope  & \pr{31}  & \nope \\
    \labelcref{sec:align_10} & Alignment 10 & 12d Solutions & \nope  & \pr{38}  & \nope \\
    \labelcref{sec:align_11} & Alignment 11 & 12d Solutions & \jaja  & Geometry Gym & \nope \\
    \labelcref{sec:align_12} & Alignment 12 & Aplitop & \jaja  & project team  & \nope \\
    \labelcref{sec:align_13} & Alignment 13 & TUM & \nope  & \pr{39}  & \nope \\
    \labelcref{sec:align_14} & Alignment 14 & project team & \jaja  & Geometry Gym  & \jaja \\
    \labelcref{sec:align_15} & Alignment 15 & Aplitop & \jaja  & -  & \nope \\
    \labelcref{sec:align_16} & Alignment 16 & Aplitop & \nope  & \pr{96}  & \nope \\
    \labelcref{sec:align_17} & Alignment 17 & 12d Solutions & \nope  & \pr{103}  & \nope \\
    \labelcref{sec:align_18} & Alignment 18 & 12d Solutions & \nope  & \pr{107}  & \nope \\
    \labelcref{sec:align_19} & Alignment 19 & ACCA & \jaja  & Geometry Gym  & \nope \\
    \labelcref{sec:align_20} & Alignment 20 & ACCA & \jaja  & Geometry Gym  & \nope \\
    \labelcref{sec:align_21} & Alignment 21 & TUM \& AutoDesk & \nope  & \pr{143}  & \nope \\
  \midrule
    \labelcref{sec:linear_1} & Linear Pos. 1 & ACCA & \jaja  & -  & \nope \\
    \labelcref{sec:linear_2} & Linear Pos. 2 & ACCA & \jaja  & -  & \nope \\
    \labelcref{sec:linear_3} & Linear Pos. 3 & ACCA & \jaja  & -  & \nope \\
    \labelcref{sec:linear_4} & Linear Pos. 4 & ACCA & \jaja  & -  & \nope \\
  \midrule
    \labelcref{sec:spatial_1} & Spatial str. 1 & project team & \jaja  & FZK Viewer & \nope \\
    \labelcref{sec:spatial_2} & Spatial str. 2 & project team & \jaja  & FZK Viewer & \jaja \\
    \labelcref{sec:spatial_3} & Spatial str. 3 & project team & \jaja  & IFCTreeView & \jaja \\
    \labelcref{sec:spatial_4} & Spatial str. 4 & project team & \jaja  & IFCTreeView & \jaja \\
    \labelcref{sec:spatial_5} & Spatial str. 5 & project team & \jaja  & FZK Viewer & \nope \\
    \labelcref{sec:spatial_6} & Spatial str. 6 & project team & \jaja  & FZK Viewer & \jaja \\
    \labelcref{sec:spatial_7} & Spatial str. 7 & project team & \jaja  & \pr{116} & \jaja \\
    \labelcref{sec:spatial_8} & Spatial str. 8 & TUM & \jaja  & \pr{65} & \nope \\
  \midrule
    \labelcref{sec:swept_1} & Swept Geom. 1 & project team & \jaja  & \pr{40} & \nope \\
    \labelcref{sec:swept_2} & Swept Geom. 2 & project team & \jaja  & - & \jaja \\
    \labelcref{sec:swept_3} & Swept Geom. 3 & Aplitop & \jaja  & \pr{106} & \nope \\
    \labelcref{sec:swept_4} & Swept Geom. 4 & ACCA & \jaja  & Geometry Gym  & \nope \\
    \labelcref{sec:swept_5} & Swept Geom. 5 & Allplan & \nope  & \pr{51} & \nope \\
    \labelcref{sec:swept_6} & Swept Geom. 6 & ACCA & \jaja  & \pr{129} & \nope \\
    \labelcref{sec:swept_7} & Swept Geom. 7 & ACCA & \nope  & \pr{132} & \nope \\
    \labelcref{sec:swept_8} & Swept Geom. 8 & ACCA & \nope  & \pr{133} & \nope \\
    \labelcref{sec:swept_9} & Swept Geom. 9 & ACCA & \nope  & \pr{134} & \nope \\
    \labelcref{sec:swept_10} & Swept Geom. 10 & ACCA & \nope  & \pr{135} & \nope \\
    \labelcref{sec:swept_11} & Swept Geom. 11 & ACCA & \nope  & \pr{136} & \nope \\
    \labelcref{sec:swept_12} & Swept Geom. 12 & ACCA & \nope  & \pr{139} & \nope \\
    \labelcref{sec:swept_13} & Swept Geom. 13 & ACCA & \nope  & \pr{140} & \nope \\
  \midrule
    \labelcref{sec:road_1} & Road 1 & AKG Software & \jaja  & \pr{52} & \nope \\
    \labelcref{sec:road_2} & Road 2 & AKG Software & \jaja  & \pr{52} & \nope \\
    \labelcref{sec:road_3} & Road 3 & project team & \jaja  & - & \jaja \\
    \labelcref{sec:road_4} & Road 4 & project team & \jaja  & - & \jaja \\
    \labelcref{sec:road_5} & Road 5 & project team & \jaja  & - & \jaja \\
    \labelcref{sec:marine_1} & Marine 1 & project team & \nope  & \pr{94} & \jaja \\
    \labelcref{sec:marine_2} & Marine 2 & project team & \jaja  & - & \jaja \\
  \midrule
    \labelcref{sec:earth_0} & Earthworks 0 & 12d Solutions & \jaja  & TUM OIP & \nope \\
    \labelcref{sec:earth_1} & Earthworks 1 & 12d Solutions & \jaja  & TUM OIP & \nope \\
    \labelcref{sec:earth_2} & Earthworks 2 & 12d Solutions & \jaja  & TUM OIP & \nope \\
    \labelcref{sec:earth_3} & Earthworks 3 & BricSys & \jaja  & 12d Solutions & \nope \\
    \labelcref{sec:borehole_1} & Geotechnics 1 & Geometry Gym & \jaja  & - & \jaja \\
    \labelcref{sec:geological_1} & Geology 1 & Geometry Gym & \nope  & - & \jaja \\
  \midrule
    \labelcref{sec:survey_1} & Geod. survey 1 & 12d Solutions & \jaja  & TUM OIP & \nope \\
    \labelcref{sec:survey_2} & Geod. survey 2 & 12d Solutions & \jaja  & \pr{56} & \nope \\
    \labelcref{sec:survey_3} & Geod. survey 3 & 12d Solutions & \jaja  & \pr{112} & \nope \\
    \labelcref{sec:survey_4} & Geod. survey 4 & 12d Solutions & \jaja  & RDF Viewer & \nope \\
  \midrule
    \labelcref{sec:drainage_1} & Drainage 1 & project team & \jaja  & - & \jaja \\
    \labelcref{sec:drainage_2} & Drainage 2 & project team & \jaja  & - & \jaja \\
  \midrule
    \labelcref{sec:align_dtm_1} & Combination 1 & 12d Solutions & \nope  & - & \nope \\
    \labelcref{sec:align_dtm_2} & Combination 2 & Card-1 & \nope  & - & \nope \\
    \labelcref{sec:georef_earth_1} & Combination 3 & 12d Solutions & \jaja  & - & \nope \\
    \labelcref{sec:georef_earth_2} & Combination 4 & 12d Solutions & \nope  & - & \nope \\
  \midrule
    \labelcref{sec:other_1} & Cost items 1 & project team & \nope  & \pr{53} & \jaja \\
    \labelcref{sec:other_2} & Properties 2 & ACCA & \nope  & \pr{120} & \nope \\
    \labelcref{sec:other_3} & Properties 3 & ACCA & \nope  & \pr{121} & \nope \\
  \bottomrule
\end{tabularx}

\subsection{Coverage}
\label{sec:coverage}

\Cref{tab:coverage} lists all IFC entities occurring in all accepted and approved unit tests.
The number next to each entity shows how many times this entity occurs\footnotemark{}.

\footnotetext{The content has been joined in Appendix C of the deliverable with the covered entities by the IFC Rail Phase 2 project but is also included here for completeness and automatic updates in the future.}

\begin{tabularx}{\textwidth}{rX}
  \caption{The covered entities of all unit test files combined.}\label{tab:coverage} \\
  \toprule
    Occurance & Entity \\
  \midrule
  \endfirsthead
  \caption[]{\emph{(cont.)} The covered entities of all unit test files combined.} \\
  \toprule
    Occurance & Entity \\
  \midrule
  \endhead
  \bottomrule
  \multicolumn{2}{r}{\footnotesize(continued)}
  \endfoot
  \bottomrule
  \endlastfoot
  \input{../Report/coverageTable.tex}
\end{tabularx}

\subsection{Schematics}
\label{sec:schematics}

\Cref{fig:schematics} represents the dependencies between the individual unit tests.
Unit tests are represented as boxes, while the arrows point from the dependant to the independent unit test.
Each unit test has a link to the corresponding section of this document in the upper right corner.

\tikzset{every label/.style={xshift=3ex, text width=10ex, align=left, 
                             inner sep=1pt, font=\footnotesize}}
\tikzset{pr/.style={dashed}}
\tikzset{star/.style={draw=none,anchor=north east,inner sep=2pt}}
\begin{figure}
  \begin{forest}
    for tree={              % style of tree nodes
      font=\footnotesize,
      draw, semithick, rounded corners,minimum width=7em,
            align = center,
        inner sep = 1pt,
                            % style of tree (edges, distances, direction)
             edge = {draw, semithick, latex'-},
    parent anchor = east,
     child anchor = west,
             grow = east,
    forked edge,            % for forked edge
            l sep = 6mm,   % level distance
         fork sep = 3mm,    % distance from parent to branching point
                }
    [Setup-1,name=boiler1,label=\labelcref{sec:project_setup_1}
      [Swept-3,name=swept3,label=\labelcref{sec:swept_3},pr
      ],
      [Road-2,name=road2,label=\labelcref{sec:road_2},pr
      ],
      [Road-1,name=road1,label=\labelcref{sec:road_1},pr
      ],
      [Spatial-1,name=spatial1,label=\labelcref{sec:spatial_1}
        [Spatial-8,name=spatial8,label=\labelcref{sec:spatial_8},pr
        ],
        [Spatial-7,name=spatial7,label=\labelcref{sec:spatial_7},pr
        ] { \node[star] at (.north east) {*}; },
        [Spatial-6,name=spatial6,label=\labelcref{sec:spatial_6}
        ] { \node[star] at (.north east) {*}; },
        [Spatial-5,name=spatial5,label=\labelcref{sec:spatial_5}
        ],
        [Spatial-4,name=spatial4,label=\labelcref{sec:spatial_4}
          [Road-4,name=road4,label=\labelcref{sec:road_4}
          ] { \node[star] at (.north east) {*}; },
          [Road-3,name=road3,label=\labelcref{sec:road_3}
          ] { \node[star] at (.north east) {*}; },
          [Drainage-2,name=drainage2,label=\labelcref{sec:drainage_2}
          ] { \node[star] at (.north east) {*}; },
          [Drainage-1,name=drainage1,label=\labelcref{sec:drainage_1}
          ] { \node[star] at (.north east) {*}; },
        ] { \node[star] at (.north east) {*}; },
        [Spatial-3,name=spatial3,label=\labelcref{sec:spatial_3}
          [CostItems-1,name=other1,label=\labelcref{sec:other_1},pr
          ] { \node[star] at (.north east) {*}; },
        ],
        [Spatial-2,name=spatial2,label=\labelcref{sec:spatial_2}
        ],
      ],
      [Bridge-1,name=bridge1,label=\labelcref{sec:bridge_1},pr
      ],
      [Marine-2,name=marine2,label=\labelcref{sec:marine_2}
      ] { \node[star] at (.north east) {*}; },
      [Marine-1,name=marine1,label=\labelcref{sec:marine_1},pr
      ] { \node[star] at (.north east) {*}; },
      [Properties-1,name=other2,label=\labelcref{sec:other_2}
        [Properties-2,name=other3,label=\labelcref{sec:other_3}
        ],
      ],
      [Setup-2,name=boiler2,label=\labelcref{sec:project_setup_2}
      ],
      [GeoRef-1,name=georef1,label=\labelcref{sec:georeferencing_1}
        [Survey-1,name=survey1,label=\labelcref{sec:survey_1}
          [Survey-2,name=survey2,label=\labelcref{sec:survey_2},pr
            [Survey-4,name=survey4,label=\labelcref{sec:survey_4}
            ],
            [Survey-3,name=survey3,label=\labelcref{sec:survey_3},pr
            ]{
              \draw[-latex'] (survey4.60) -- (.300);
            },
          ],
        ],
        [GeoRef-3,name=georef3,label=\labelcref{sec:georeferencing_3},pr          
          [GeoRef-4,name=georef4,label=\labelcref{sec:georeferencing_4},pr
          ]
        ],
      ] {
        \draw[-latex'] () to[out=north,in=south] (tin1);
      },
      [TIN-1,name=tin1,label=\labelcref{sec:tin_1}
        [GeoRef-TIN-1,name=georeftin1,label=\labelcref{sec:georef_tin_1} 
          [GeoRef-TIN-4,name=georeftin4,label=\labelcref{sec:georef_tin_4}
          ],
          [GeoRef-TIN-3,name=georeftin3,label=\labelcref{sec:georef_tin_3}
          ],      
          [GeoRef-TIN-2,name=georeftin2,label=\labelcref{sec:georef_tin_2}
          ],
          [Combi-2,name=combi2,label=\labelcref{sec:georef_earth_1}
            [Combi-3,name=combi3,label=\labelcref{sec:georef_earth_2},pr
            ]
          ] { \draw[-latex'] () to[out=north west,in=east] (earth1); }
        ] { \draw[-latex'] () to[out=south west,in=north east] (georef1); }
      ],
      [Earthworks-3,name=earth3,label=\labelcref{sec:earth_3},pr
      ],
      [Earthworks-0,name=earth0,label=\labelcref{sec:earth_0}
        [Earthworks-2,name=earth2,label=\labelcref{sec:earth_2}
        ],
        [Earthworks-1,name=earth1,label=\labelcref{sec:earth_1}
        ],
      ],
      [Geotechnics-1,name=geotech1,label=\labelcref{sec:borehole_1}
      ] { \node[star] at (.north east) {*}; },      
      [Alignment-17,name=align17,label=\labelcref{sec:align_17},pr  
        [Alignment-18,name=align18,label=\labelcref{sec:align_18},pr
        ],
      ],   
      [Alignment-14,name=align14,label=\labelcref{sec:align_14},pr
        [Swept-2,name=swept2,label=\labelcref{sec:swept_2},pr
        ] { \node[star] at (.north east) {*}; },
        [Road-5,name=road5,label=\labelcref{sec:road_5}
        ] { \node[star] at (.north east) {*}; },
      ] { \node[star] at (.north east) {*}; },
      [Alignment-13,name=align13,label=\labelcref{sec:align_13},pr
        [Swept-1,name=swept1,label=\labelcref{sec:swept_1},pr
        ],
      ],
      [Alignment-12,name=align12,label=\labelcref{sec:align_12},pr
        [Alignment-15,name=align15,label=\labelcref{sec:align_15},pr
        ],
        [Alignment-16,name=align16,label=\labelcref{sec:align_16},pr
        ],
      ],
      [Alignment-20,name=align20,label=\labelcref{sec:align_20}
      ],
      [Alignment-19,name=align19,label=\labelcref{sec:align_19}
      ],
      [Swept-4,name=swept4,label=\labelcref{sec:swept_4}
      ],
      [Alignment-1,name=align1,label=\labelcref{sec:align_1}
        [Alignment-9,name=align9,label=\labelcref{sec:align_9},pr
          [Alignment-11,name=align11,label=\labelcref{sec:align_11}
          ],
          [Alignment-10,name=align10,label=\labelcref{sec:align_10},pr
          ]
        ],
        [Alignment-8,name=align8,label=\labelcref{sec:align_8}
        ],
        [Alignment-6,name=align6,label=\labelcref{sec:align_6}
          [Alignment-7,name=align7,label=\labelcref{sec:align_7}
          ]
        ],
        [Alignment-4,name=align4,label=\labelcref{sec:align_4}
          [Alignment-5,name=align5,label=\labelcref{sec:align_5}
          ],
        ],
        [Alignment-2,name=align2,label=\labelcref{sec:align_2}
          [Alignment-3,name=align3,label=\labelcref{sec:align_3}
          ],
        ],
      ],
    ]  
  \end{forest}
  \caption{Schematic representation of unit tests and their dependencies.
           Dashed boxes mean that they are still under review.
           A star in the upper right corner means that the unit test has been derived from a storyline.}
  \label{fig:schematics}
\end{figure}

\subsection{Overview over the Content}
\label{sec:content}
\renewcommand\contentsname{}
\tableofcontents
\clearpage

% -------------------------------------------------------------------------------
\section{Project Setup}
\label{sec:project_setup}

\subsection{Project Setup 1}
\label{sec:project_setup_1}
\markdownInput{../ProjectSetup-1/readme.md}
\clearpage

\subsection{Project Setup 2}
\label{sec:project_setup_2}
\markdownInput{../ProjectSetup-2/readme.md}
\clearpage

% -------------------------------------------------------------------------------
\section{Georeferencing and Digital Terrain Models (DTMs)}
\label{sec:georeferencing_and_dtms}

\subsection{TIN 1}
\label{sec:tin_1}
\markdownInput{../Tin-1/readme.md}
\clearpage

\subsection{Georeferencing 1}
\label{sec:georeferencing_1}
\markdownInput{../Georeferencing-1/readme.md}
\clearpage

\subsection{Georeferencing 2} % PR#85
\label{sec:georeferencing_3}
Under review.%\markdownInput{../Georeferencing-3/readme.md}
\clearpage

\subsection{Georeferencing 3} % PR#85
\label{sec:georeferencing_4}
Under review.%\markdownInput{../Georeferencing-4/readme.md}
\clearpage

\subsection{Georeferencing and DTM 1}
\label{sec:georef_tin_1}
\markdownInput{../Georeferencing-Tin-12d-1/readme.md}
\clearpage

\subsection{Georeferencing and DTM 2}
\label{sec:georef_tin_2}
\markdownInput{../Georeferencing-Tin-12d-2/readme.md}
\clearpage

\subsection{Georeferencing and DTM 3} % PR#20
\label{sec:georef_tin_3}
\markdownInput{../Georeferencing-Tin-12d-3/readme.md}
\clearpage

\subsection{Georeferencing and DTM 4} % PR#21
\label{sec:georef_tin_4}
\markdownInput{../Georeferencing-Tin-12d-4/readme.md}
\clearpage

% -------------------------------------------------------------------------------
\section{Alignment}
\label{sec:alignment}

\subsection{Alignment 1}
\label{sec:align_1}
\markdownInput{../Alignment-12d-1/readme.md}
\clearpage

\subsection{Alignment 2}
\label{sec:align_2}
\markdownInput{../Alignment-12d-2/readme.md}
\clearpage

\subsection{Alignment 3}
\label{sec:align_3}
\markdownInput{../Alignment-12d-3/readme.md}
\clearpage

\subsection{Alignment 4}
\label{sec:align_4}
\markdownInput{../Alignment-12d-4/readme.md}
\clearpage

\subsection{Alignment 5}
\label{sec:align_5}
\markdownInput{../Alignment-12d-5/readme.md}
\clearpage

\subsection{Alignment 6} % PR#25
\label{sec:align_6}
\markdownInput{../Alignment-12d-6/readme.md}
\clearpage

\subsection{Alignment 7} % PR#26
\label{sec:align_7}
\markdownInput{../Alignment-12d-7/readme.md}
\clearpage

\subsection{Alignment 8} % PR#29
\label{sec:align_8}
\markdownInput{../Alignment-12d-8/readme.md}
\clearpage

\subsection{Alignment 9} % PR#31
\label{sec:align_9}
Under review.%\markdownInput{../Alignment-12d-9/readme.md}
\clearpage

\subsection{Alignment 10} % PR#38
\label{sec:align_10}
Under review.%\markdownInput{../Alignment-12d-10/readme.md}
\clearpage

\subsection{Alignment 11} % PR#42
\label{sec:align_11}
\markdownInput{../Alignment-12d-11/readme.md}
\clearpage

\subsection{Alignment 12} % PR#96
\label{sec:align_12}
\markdownInput{../Alignment-Aplitop-1/readme.md}
\clearpage

\subsection{Alignment 13} % PR#39
\label{sec:align_13}
Under review.%\markdownInput{../UT-Alignment-3/readme.md}
\clearpage

\subsection{Alignment 14} % PR#71 *
\label{sec:align_14}
\markdownInput{../Alignment-INDOT/readme.md}
\clearpage

\subsection{Alignment 15} % PR#96
\label{sec:align_15}
\markdownInput{../Alignment-Aplitop-2/readme.md}
\clearpage

\subsection{Alignment 16} % PR#96
\label{sec:align_16}
Under review.%\markdownInput{../Alignment-Aplitop-3/readme.md}
\clearpage

\subsection{Alignment 17} % PR#103
\label{sec:align_17}
Under review.%\markdownInput{../Georeferencing-Alignment-12d-1/readme.md}
\clearpage

\subsection{Alignment 18} % PR#107
\label{sec:align_18}
Under review.%\markdownInput{../Georeferencing-Alignment-12d-2/readme.md}
\clearpage

\subsection{Alignment 19} % PR#122
\label{sec:align_19}
\markdownInput{../Alignment-ACCA-1/README.md}
\clearpage

\subsection{Alignment 20} % PR#123
\label{sec:align_20}
\markdownInput{../Alignment-ACCA-2/README.md}
\clearpage

\subsection{Alignment 21} % PR#143
\label{sec:align_21}
Under review.%\markdownInput{../Alignment-TUMAdsk-1/readme.md}
\clearpage

% -------------------------------------------------------------------------------
\section{Linear Placement}
\label{sec:linear_placement}

\subsection{Linear Placement 1} 
\label{sec:linear_1}
\markdownInput{../LinearPlacement-1/README.md}
\clearpage

\subsection{Linear Placement 2} 
\label{sec:linear_2}
\markdownInput{../LinearPlacement-2/README.md}
\clearpage

\subsection{Linear Placement 3} 
\label{sec:linear_3}
\markdownInput{../LinearPlacement-3/README.md}
\clearpage

\subsection{Linear Placement 4} 
\label{sec:linear_4}
\markdownInput{../LinearPlacement-4/README.md}
\clearpage

% -------------------------------------------------------------------------------
\section{Spatial Structure}
\label{sec:spatial_structure}

\subsection{Spatial Structure 1}
\label{sec:spatial_1}
\markdownInput{../SpatialStructure-1/readme.md}
\clearpage

\subsection{Spatial Structure 2}
\label{sec:spatial_2}
\markdownInput{../SpatialStructure-2/readme.md}
\clearpage

\subsection{Spatial Structure 3}
\label{sec:spatial_3}
\markdownInput{../SpatialStructure-3/readme.md}
\clearpage

\subsection{Spatial Structure 4}
\label{sec:spatial_4}
\markdownInput{../SpatialStructure-4/readme.md}
\clearpage

\subsection{Spatial Structure 5}
\label{sec:spatial_5}
\markdownInput{../SpatialStructure-5/readme.md}
\clearpage

\subsection{Spatial Structure 6}
\label{sec:spatial_6}
\markdownInput{../SpatialStructure-6/readme.md}
\clearpage

\subsection{Spatial Structure 7} % PR#116 *
\label{sec:spatial_7}
Under review.%\markdownInput{../SpatialStructure-7/README.md}
\clearpage

\subsection{Spatial Structure 8} % PR#65
\label{sec:spatial_8}
Under review.%\markdownInput{../CrossDisciplinary-1/README.md}
\clearpage

% -------------------------------------------------------------------------------
\section{Swept Geometry}
\label{sec:swept_geometry}


\subsection{Swept Geometry 1} % PR#40
\label{sec:swept_1}
Under review.%\markdownInput{../UT-Road-1/readme.md}
\clearpage

\subsection{Swept Geometry 2} % PR#71 *
\label{sec:swept_2}
\markdownInput{../SectionedSurface-INDOT/readme.md}
\clearpage

\subsection{Swept Geometry 3} % PR#106 *
\label{sec:swept_3}
Under review.%\markdownInput{../SectionedSurface-Aplitop-1/readme.md}
\clearpage

\subsection{Swept Geometry 4} % PR #128
\label{sec:swept_4}
\markdownInput{../FixedReferenceSweptAreaSolid-1/README.md}
\clearpage

\subsection{Swept Geometry 5} % PR#51
\label{sec:swept_5}
Under review.%\markdownInput{../SectionedSolidHorizontal-1/readme.md}
\clearpage

\subsection{Swept Geometry 6} % PR #129
\label{sec:swept_6}
Under review.%\markdownInput{../SectionedSolidHorizontal-1/README.md}
\clearpage

\subsection{Swept Geometry 7} % PR #132
\label{sec:swept_7}
Under review.%\markdownInput{../SectionedSolidHorizontal-2/README.md}
\clearpage

\subsection{Swept Geometry 8} % PR #133
\label{sec:swept_8}
Under review.%\markdownInput{../SectionedSolidHorizontal-3/README.md}
\clearpage

\subsection{Swept Geometry 9} % PR #134
\label{sec:swept_9}
Under review.%\markdownInput{../SectionedSolidHorizontal-4/README.md}
\clearpage

\subsection{Swept Geometry 10} % PR #135
\label{sec:swept_10}
Under review.%\markdownInput{../SectionedSolidHorizontal-5/README.md}
\clearpage

\subsection{Swept Geometry 11} % PR #136
\label{sec:swept_11}
Under review.%\markdownInput{../SectionedSolidHorizontal-6/README.md}
\clearpage

\subsection{Swept Geometry 12} % PR #139
\label{sec:swept_12}
Under review.%\markdownInput{../DirectrixDerivedReferenceSweptAreaSolid-1/README.md}
\clearpage

\subsection{Swept Geometry 13} % PR #140
\label{sec:swept_13}
Under review.%\markdownInput{../DirectrixDerivedReferenceSweptAreaSolid-2/README.md}
\clearpage

% -------------------------------------------------------------------------------
\section{Domain Specifics}
\label{sec:domain_specifics}

\subsection{Road 1}  % PR#52
\label{sec:road_1}
Under review.%\markdownInput{../UT-IfcCourse/readme.md}
\clearpage

\subsection{Road 2} % PR#52
\label{sec:road_2}
Under review.%\markdownInput{../UT-IfcPavement/readme.md}
\clearpage

\subsection{Road 3} % *
\label{sec:road_3}
\markdownInput{../MCON-1/readme.md}
\clearpage

\subsection{Road 4} % *
\label{sec:road_4}
\markdownInput{../MCON-2/readme.md}
\clearpage

\subsection{Road 5} % PR#88 *
\label{sec:road_5}
\markdownInput{../RumbleStrip-INDOT/readme.md}
\clearpage

\subsection{Marine 1} % PR#94
\label{sec:marine_1}
Under review.%\markdownInput{../MarineFurniture-1/readme.md}
\clearpage

\subsection{Marine 2} % PR#97
\label{sec:marine_2}
\markdownInput{../MarineFurniture-2/readme.md}
\clearpage

% -------------------------------------------------------------------------------
\section{Earthworks, Geology \& Geotechnics}
\label{sec:earthworks}

\subsection{Earthworks 0}
\label{sec:earth_0}
\markdownInput{../Earthworks-12d-0/readme.md}
\clearpage

\subsection{Earthworks 1}
\label{sec:earth_1}
\markdownInput{../Earthworks-12d-1/readme.md}
\clearpage

\subsection{Earthworks 2}
\label{sec:earth_2}
\markdownInput{../Earthworks-12d-2/readme.md}
\clearpage

\subsection{Earthworks 3}
\label{sec:earth_3}
\markdownInput{../Earthworks-01/readme.md}
\clearpage

\subsection{Borehole 1}
\label{sec:borehole_1}
\markdownInput{../Borehole-1/readme.md}
\clearpage

\subsection{Geological 1} % PR#141
\label{sec:geological_1}
Under review.%\markdownInput{../Geological-1/readme.md}
\clearpage

% -------------------------------------------------------------------------------
\section{Geodetic Survey}
\label{sec:geodetic_survey}

\subsection{Survey 1}
\label{sec:survey_1}
\markdownInput{../Georeferencing-Survey-12d-1/readme.md}
\clearpage

\subsection{Survey 2} % PR#56
\label{sec:survey_2}
Under review.%\markdownInput{../Georeferencing-Survey-12d-2/readme.md}
\clearpage

\subsection{Survey 3} % PR#112
\label{sec:survey_3}
Under review.%\markdownInput{../Georeferencing-Survey-12d-3/readme.md}
\clearpage

\subsection{Survey 4} % PR#118
\label{sec:survey_4}
\markdownInput{../Georeferencing-Survey-12d-4/readme.md}
\clearpage


% -------------------------------------------------------------------------------
\section{Drainage System}
\label{sec:drainage_system}

\subsection{Drainage System 1}
\label{sec:drainage_1}
\markdownInput{../DrainageSystem-1/Readme.md}
\clearpage

\subsection{Drainage System 2}
\label{sec:drainage_2}
\markdownInput{../DrainageSystem-2/Readme.md}
\clearpage

% -------------------------------------------------------------------------------
\section{Combinations}
\label{sec:combinations}

\subsection{Combination 1: Georeferencing and Alignment and DTM} % PR#43
\label{sec:align_dtm_1}
Under review.%\markdownInput{../Georeferencing-Alignment-Tin-12d-1/readme.md}
\clearpage

\subsection{Combination 2: Alignment and DTM} % PR#59
\label{sec:align_dtm_2}
Under review.%\markdownInput{../Alignment-Tin-1/readme.md}
\clearpage

\subsection{Combination 3: Georeferencing and Earthworks}
\label{sec:georef_earth_1}
\markdownInput{../Georeferencing-Earthworks-12d-1/readme.md}
\clearpage

\subsection{Combination 4: Georeferencing and Earthworks 2} % PR#50
\label{sec:georef_earth_2}
Under review.%\markdownInput{../Georeferencing-Earthworks-12d-2/readme.md}
\clearpage

% -------------------------------------------------------------------------------
\section{Other}
\label{sec:other}

\subsection{Other 1: Cost items} % PR#53 *
\label{sec:other_1}
Under review.%\markdownInput{../UT-CostItem-COWI/readme.md}
\clearpage

\subsection{Other 2: Properties 1} % PR#120
\label{sec:other_2}
\markdownInput{../Properties-1/README.md}
\clearpage

\subsection{Other 3: Properties 2} % PR#121
\label{sec:other_3}
\markdownInput{../Properties-2/README.md}
\clearpage


\end{document}
